% GENERATED FILE. Edit canonical lessons, then run npm run generate:books.
% canonical-source-hash: fnv1a64:1b3ebfb318d576fb
% canonical-lessons: FR-C395-utilisateur, FR-C395-mot-de-passe, FR-C395-identifiant, FR-C395-connecter, FR-C395-acceder

\chapter{Mon compte en ligne}
\label{ch:fr-mon-compte-en-ligne}
\hlchaptermodality{395}

\begin{chapteropening}
\textbf{By the end of this chapter:} \emph{I can manage an online account: the user, the password, the username, logging in and getting access.}

Complete the last lesson of chapter 395: I can manage an online account: the user, the password, the username, logging in and getting access.
\end{chapteropening}

\section[l'utilisateur]{l'utilisateur — the user}

\label{lesson:FR-C395-utilisateur}



\begin{quote}
Before the new one: say the French for to pay back, then the French for the loan.
\end{quote}

\subsection*{You'll want to know: l'utilisateur}
\textbf{l'utilisateur} — \textquotedblleft{}the user\textquotedblright{}.

\begin{quote}
Chaque utilisateur a son propre profil.
\end{quote}

\begin{quote}
\emph{Each user has their own profile.}
\end{quote}

Feminine \emph{utilisatrice}.

\begin{cousinweb}
\emph{utilisateur} is built on \emph{utiliser} \textquotedblleft{}to use\textquotedblright{}, from \emph{utile} \textquotedblleft{}useful\textquotedblright{}, which comes from Latin \emph{ūtilis} \textquotedblleft{}useful\textquotedblright{}, from \emph{ūtī} \textquotedblleft{}to use\textquotedblright{} (Lewis \& Short).
\end{cousinweb}

\begin{grammarlens}[title={What you've built}]
One more word for telling, explaining and describing at B1.
\end{grammarlens}

\subsection*{Your turn}
\begin{itemize}
\raggedright
  \item \emph{Say it:} \emph{l'utilisateur}
  \item \emph{Say it:} \emph{l'utilisateur}, once more
  \item \emph{Say it:} a sentence of your own with \emph{l'utilisateur}
  \item \emph{Recall:} say the French for to pay back, then the French for the loan, then say \emph{l'utilisateur} again
\end{itemize}

\begin{tcolorbox}[breakable,colback=teal!4,colframe=teal!35!black,title={Before you move on}]
What does l'utilisateur mean? (\textquotedblleft{}The user\textquotedblright{}.) Say it once more.
\end{tcolorbox}
\section[le mot de passe]{le mot de passe — the password}

\label{lesson:FR-C395-mot-de-passe}



\begin{quote}
Before the new one: say the French for the loan, then the French for the user.
\end{quote}

\subsection*{You'll want to know: le mot de passe}
\textbf{le mot de passe} — \textquotedblleft{}the password\textquotedblright{}.

\begin{quote}
J'ai encore oublié mon mot de passe.
\end{quote}

\begin{quote}
\emph{I forgot my password again.}
\end{quote}

The plural is \emph{les mots de passe}.

\begin{cousinweb}
\emph{mot de passe} is a word that lets you pass. \emph{mot} comes from Late Latin \emph{muttum} \textquotedblleft{}a mutter, a grunt\textquotedblright{} (EtymDB; Lewis \& Short), and \emph{passer} from Vulgar Latin \emph{passāre}, built on \emph{passus} \textquotedblleft{}a step\textquotedblright{} (Lewis \& Short).
\end{cousinweb}

\begin{grammarlens}[title={What you've built}]
One more word for telling, explaining and describing at B1.
\end{grammarlens}

\subsection*{Your turn}
\begin{itemize}
\raggedright
  \item \emph{Say it:} \emph{le mot de passe}
  \item \emph{Say it:} \emph{le mot de passe}, once more
  \item \emph{Say it:} a sentence of your own with \emph{le mot de passe}
  \item \emph{Recall:} say the French for the loan, then the French for the user, then say \emph{le mot de passe} again
\end{itemize}

\begin{tcolorbox}[breakable,colback=teal!4,colframe=teal!35!black,title={Before you move on}]
What does le mot de passe mean? (\textquotedblleft{}The password\textquotedblright{}.) Say it once more.
\end{tcolorbox}
\section[l'identifiant]{l'identifiant — the username, the login}

\label{lesson:FR-C395-identifiant}



\begin{quote}
Before the new one: say the French for the user, then the French for the password.
\end{quote}

\subsection*{You'll want to know: l'identifiant}
\textbf{l'identifiant} — \textquotedblleft{}the username, the login\textquotedblright{}.

\begin{quote}
Chaque employé a un identifiant et un mot de passe.
\end{quote}

\begin{quote}
\emph{Each employee has a username and a password.}
\end{quote}

Masculine. From \emph{identifier}, to identify.

\begin{cousinweb}
\emph{identifiant} is the present participle of \emph{identifier} used as a noun. The verb comes from Medieval Latin \emph{identificāre}, built on Latin \emph{īdem} \textquotedblleft{}the same\textquotedblright{} (Lewis \& Short).
\end{cousinweb}

\begin{grammarlens}[title={What you've built}]
One more word for telling, explaining and describing at B1.
\end{grammarlens}

\subsection*{Your turn}
\begin{itemize}
\raggedright
  \item \emph{Say it:} \emph{l'identifiant}
  \item \emph{Say it:} \emph{l'identifiant}, once more
  \item \emph{Say it:} a sentence of your own with \emph{l'identifiant}
  \item \emph{Recall:} say the French for the user, then the French for the password, then say \emph{l'identifiant} again
\end{itemize}

\begin{tcolorbox}[breakable,colback=teal!4,colframe=teal!35!black,title={Before you move on}]
What does l'identifiant mean? (\textquotedblleft{}The username, the login\textquotedblright{}.) Say it once more.
\end{tcolorbox}
\section[connecter]{connecter — to connect; to log in (reflexive, se connecter)}

\label{lesson:FR-C395-connecter}



\begin{quote}
Before the new one: say the French for the password, then the French for the username.
\end{quote}

\subsection*{You'll want to know: connecter}
\textbf{connecter} — \textquotedblleft{}to connect; to log in (reflexive, se connecter)\textquotedblright{}.

\begin{quote}
Je me connecte à mon compte tous les matins.
\end{quote}

\begin{quote}
\emph{I log in to my account every morning.}
\end{quote}

\emph{Se connecter} — to log in; \emph{se déconnecter} — to log out.

\begin{cousinweb}
\emph{connecter} is built on Latin \emph{cōnectere} \textquotedblleft{}to bind together\textquotedblright{} (Lewis \& Short), from \emph{nectere} \textquotedblleft{}to tie\textquotedblright{}, partly through English \emph{connect}.
\end{cousinweb}

\begin{grammarlens}[title={What you've built}]
One more word for telling, explaining and describing at B1.
\end{grammarlens}

\subsection*{Your turn}
\begin{itemize}
\raggedright
  \item \emph{Say it:} \emph{connecter}
  \item \emph{Say it:} \emph{connecter}, once more
  \item \emph{Say it:} a sentence of your own with \emph{connecter}
  \item \emph{Recall:} say the French for the password, then the French for the username, then say \emph{connecter} again
\end{itemize}

\begin{tcolorbox}[breakable,colback=teal!4,colframe=teal!35!black,title={Before you move on}]
What does connecter mean? (\textquotedblleft{}To connect; to log in (reflexive, se connecter)\textquotedblright{}.) Say it once more.
\end{tcolorbox}
\section[accéder]{accéder — to access, to reach}

\label{lesson:FR-C395-acceder}



\begin{quote}
Before the new one: say the French for the username, then the French for to connect.
\end{quote}

\subsection*{You'll want to know: accéder}
\textbf{accéder} — \textquotedblleft{}to access, to reach\textquotedblright{}.

\begin{quote}
Sans mot de passe, on ne peut pas accéder au site.
\end{quote}

\begin{quote}
\emph{Without a password, you cannot get into the website.}
\end{quote}

\emph{Accéder à}, with \emph{à}. \emph{J'accède}.

\begin{cousinweb}
\emph{accéder} is a learned borrowing from Latin \emph{accēdere} \textquotedblleft{}to go toward, to approach\textquotedblright{} (EtymDB; Lewis \& Short), built on \emph{cēdere} \textquotedblleft{}to go\textquotedblright{}.
\end{cousinweb}

\begin{grammarlens}[title={What you've built}]
One more word for telling, explaining and describing at B1.
\end{grammarlens}

\subsection*{Your turn}
\begin{itemize}
\raggedright
  \item \emph{Say it:} \emph{accéder}
  \item \emph{Say it:} \emph{accéder}, once more
  \item \emph{Say it:} a sentence of your own with \emph{accéder}
  \item \emph{Recall:} say the French for the username, then the French for to connect, then say \emph{accéder} again
\end{itemize}

\begin{tcolorbox}[breakable,colback=teal!4,colframe=teal!35!black,title={Before you move on}]
What does accéder mean? (\textquotedblleft{}To access, to reach\textquotedblright{}.) Say it once more.
\end{tcolorbox}
